\section{ENABLING TECHNOLOGIES, APPLICATION REQUIREMENTS, AND 6G NETWORK EVIDENCE}
\label{sec:technologies_applications_6g}
Section~\ref{sec:validation_reproducibility} established how far the reported
systems were tested. This section follows what those systems enable. It connects
the mechanisms that create useful observables to the requirements of their
application domains and then examines how far the evidence reaches toward a 6G
network role. A study may contribute to several technology and application
labels, so the counts describe coverage within the corpus.

The synthesis first explains how waveforms, optical components, and spatial
control shape the observable. It then follows integration through devices,
infrastructure, and inference. The application map translates these mechanisms
into operating conditions, and the final subsection explains what those labels
mean for 6G network evidence.

% ============================================================================
% PLANNED FIGURE 8 --- FROM OPTICAL OBSERVABLES TO NETWORK EVIDENCE
% DO NOT DRAW IN THIS VERSION
% Stable label fig:technology_application_chain
% Blueprint ID FIG-OISAC-07
% Placement after the section framing and before the detailed synthesis
% Data authority Phase F s6_enabling_technologies.csv,
% s6_application_domains.csv, and s7_six_g_relevance.csv. Technology and
% application counts are multilabel. The 6G relevance gate is exclusive over
% all 206 studies.
%
% Purpose
% Give the reader one cross layer map showing how each layer contributes to a
% joint service. Marginal multilabel counts describe layer coverage.
%
% Layout
% Use six layers from left to right with uniform technical progression arrows
% and dashed reverse requirement arrows.
% 1. Generation and transport with photonic THz 68, coherent optics 64, photonic
%    integration 20, fiber DAS/infrastructure reuse 22, other 19.
% 2. Waveform and physical observables with FMCW or chirped 66, OFDM or
%    multicarrier 56, and phase, polarization, backscatter, spatial, power, and
%    image observables.
% 3. Spatial control with beamforming 13, OPA 8, MIMO 7, and RIS or ORIS 1.
% 4. Inference and system models with ML or AI 20 and digital twin 2. Show
%    calibration, training distribution, uncertainty, computation, and latency.
% 5. Four application requirement families covering network and mobility,
%    infrastructure and environment, positioning and human centered services,
%    and extreme settings. Table VII retains the thirteen individual domains.
% 6. The exclusive 6G evidence gate with direct 138, inferential 64, weak 1,
%    and not applicable 3. Follow it with network role, timing, orchestration,
%    interoperability, paired function evidence, and validated service.
%    The closing label reads ``6G relevance marks corpus positioning.
%    Conformance, interoperability, and readiness require further evidence.''
%
% Activation rule
% After the figure is produced and audited, replace the count inventories at
% the start of subsections VII-A, VII-B, and VII-D with the following lead in.
% Figure~\ref{fig:technology_application_chain} organizes the section as a
% cross layer evidence chain from generation and observable formation to
% application requirements and a bounded 6G relevance gate.
%
% Future caption
% Cross layer organization of O-ISAC enabling technologies, observables,
% application requirements, and 6G evidence. Generation and transport create or
% preserve observables. Waveform design, spatial control, and inference shape
% their use. Application requirements return constraints to the system. The
% labels describe corpus coverage, and the arrows show technical progression.
% The final gate reports direct, inferential, weak, and not applicable relevance
% as 138, 64, 1, and 3 studies. Conformance, interoperability, paired operation,
% and deployment require additional evidence.
%
% Visual constraints
% Use a conceptual chain instead of a Sankey diagram and keep every arrow width
% uniform. Connect technologies at the layer level because the current data
% authority supplies marginal counts. Keep wireless propagation explicit after
% photonic carrier generation. Present 6G relevance as corpus positioning. Use
% grayscale safe shapes and direct layer labels.
% ============================================================================

\subsection{Technologies That Create and Shape Observables}

The most frequent technology labels occupy complementary points in the signal
chain. Photonic terahertz and microwave photonic generation appear in 68
studies. Frequency modulated continuous wave and related chirped processing
appear in 66, coherent optical techniques in 64, and OFDM or related
multicarrier methods in 56. Photonic generation creates or translates the
carrier. Chirped processing maps delay or motion to beat frequency. Coherent
reception preserves phase, and OFDM provides controllable time and frequency
resources. A single architecture can combine several of these mechanisms.

OFDM supports joint operation by placing payload and sensing resources within
one signal structure. The sensing observable may come from pilots, payload
tones, cyclic prefixes, reserved subcarriers, or a dedicated processing branch.
This choice determines which samples are known at the receiver and how
interference enters the estimate. Optical modulation constraints and the
measurement plane used for SNR shape the interpretation. Embedded pilot fiber
sensing and indoor optical positioning show how the same multicarrier principle
serves different observables and geometries
\cite{OISAC_SCR00007,OISAC_SCR01024}.

Chirped and coherent systems draw on phase continuity and a shared timing
reference. A returned chirp can reveal range or motion as the transmitted signal
carries data. A coherent front end preserves small phase changes for distributed
sensing and demodulation. Photonic generation can extend the reference to
wideband millimeter wave or terahertz carriers. Full photonic dechirping and
coherent fiber vibration sensing illustrate these two uses
\cite{OISAC_SCR01000,OISAC_SCR00647}. Their performance depends on laser
stability, reference distribution, conversion efficiency, receiver bandwidth,
and calibration.

The optical stage governs carrier generation and reference stability. Once the
carrier is radiated, the communication and sensing paths follow microwave,
millimeter wave, or terahertz propagation. In a free space optical link, the
optical beam itself traverses the channel. System evaluation should therefore
follow the signal across both physical domains and identify where power, SNR,
bandwidth, and timing are measured. This view makes the interface between
components part of the joint system.

Spatial control changes both the communication link budget and the sensing
geometry. Beamforming appears in 13 studies, optical phased arrays in 8,
multiple input multiple output systems in 7, and a reconfigurable optical
intelligent surface in 1. These labels describe different functions.
Beamforming directs energy. An optical phased array provides one steering
mechanism. Multiple input multiple output processing uses several spatial
dimensions, and an intelligent surface reshapes an intermediate propagation
boundary. A system may combine these roles.

A narrow beam improves received power and angular selectivity while making
acquisition, tracking, blockage, and calibration central to joint performance.
The joint use of an optical phased array and OFDM and the use of frequency comb
beamforming connect steering with ranging or imaging
\cite{OISAC_SCR00941,OISAC_SCR00362}. Multiple input multiple output processing
adds communication streams, sensing diversity, or separation among targets and
paths. Structured fields and hybrid systems that combine free space
optics with fiber Bragg gratings place these spatial dimensions at different
points in the architecture \cite{OISAC_SCR00268,OISAC_SCR00277}. Meaningful
comparison therefore requires the independent dimensions, geometry, and
calibration method.

The intelligent surface evidence consists of one study that models a mobile
free space optical path controlled by an optical intelligent surface
\cite{OISAC_SCR00914}. A separate laboratory demonstration uses mechanically
rotated Risley prisms to steer a photonically generated terahertz signal
\cite{OISAC_SCR00953}. Both mechanisms redirect energy, and each uses different
surface physics, propagation laws, control variables, and validation settings.
Their separate treatment preserves the physical meaning and provides two
starting points for broader experimental comparison.

\subsection{Integration Beyond the Waveform}

Integration and inference move joint design beyond waveform selection.
Photonic integration appears in 20 studies, fiber distributed acoustic sensing
in 22, and machine learning or artificial intelligence in 20. Two studies use
digital twins. A further 19 studies contain specialized mechanisms grouped as
other because they fall outside the broad technology labels. Together, these
categories cover device integration, infrastructure reuse, inference, and a
diverse long tail.

Photonic integration can place generation, reference distribution, beam
control, and detection on a compact hardware path. This arrangement reduces
duplicated paths and brings packaging, thermal drift, crosstalk, and calibration
into the same design problem. Microcomb synchronization coordinates wireless
data with radar imaging. A distributed W band system uses radar estimates for
beam alignment and flexible data service
\cite{OISAC_SCR00406,OISAC_SCR00502}. Both examples show how a shared reference
can connect communication and sensing hardware.

Fiber distributed acoustic sensing follows a complementary route by using the
transmission medium as a spatial sensor. Phase, backscatter, and polarization
can reveal events along the route as the fiber carries traffic. Coherent
multicore transmission with distributed sensing and operational submarine cable
monitoring illustrate two scales of this infrastructure reuse
\cite{OISAC_SCR00346,OISAC_SCR00957}. The central design question becomes how to
preserve the optical observable together with traffic continuity.

Learning turns optical observations into event labels, positions, field
estimates, channel predictions, and resource decisions. Each task is best read
with its definition, reference data, training distribution, evaluation shift,
uncertainty, and inference cost. Fiber diagnosis, federated visible light
positioning, and wearable optical sensing show the range from distributed
infrastructure to human centered use
\cite{OISAC_SCR01006,OISAC_SCR01081,OISAC_SCR00376}.

A neural network can also take part in both functions. In one coherent link,
the model compensates transmission impairments and estimates distributed power
and abnormal loss \cite{OISAC_SCR00949}. The two outputs share the link while
retaining separate targets and reference data. Digital twins add another model
layer. One represents vehicular edge resources, and another provides a model
aided physical reference for coherent fiber sensing
\cite{OISAC_SCR01059,OISAC_SCR01185}. Clear reporting of the represented state,
update process, mismatch model, and resulting decision gives these models an
explicit system role.

\subsection{Applications as Operating Requirements}

Applications translate the mechanisms above into conditions that a joint
service must satisfy. Table~\ref{tab:application_requirements} retains all
thirteen application domains and maps each one to its joint role and defining
evaluation conditions. A fiber platform for vehicles, for
example, carries high speed data and measures distributed temperature
\cite{OISAC_SCR00123}. Human centered devices connect communication or remote
alerts with motion, falls, air pressure, physiology, or gesture
\cite{OISAC_SCR00451,OISAC_SCR00294,OISAC_SCR00376}. Each application therefore
changes what the experiment must demonstrate, even when both functions share
the same optical interface.

\begin{table*}[!t]
\caption{Application domains and operating requirements in the reviewed corpus}
\label{tab:application_requirements}
\centering
\footnotesize
\setlength{\tabcolsep}{3pt}
\renewcommand{\arraystretch}{1.05}
\begin{tabularx}{\textwidth}{@{}
>{\raggedright\arraybackslash}p{0.16\textwidth}
>{\raggedright\arraybackslash}p{0.39\textwidth}
>{\raggedright\arraybackslash}X
>{\centering\arraybackslash}p{0.095\textwidth}@{}}
\toprule
\multicolumn{4}{@{}l}{\emph{Network and mobility}} \\
\addlinespace[1pt]
\textbf{Domain, n (\%)} & \textbf{Common communication and sensing role} &
\textbf{Key operating condition} & \textbf{Studies} \\
\midrule
\textbf{6G and access} 100 (48.5\%) & Access and transport with channel, target,
position, or environment sensing & Timing, orchestration, continuity, and
evidence for both functions & \cite{OISAC_SCR00277,OISAC_SCR01049} \\
\addlinespace[1.25pt]
\textbf{Vehicular} 41 (19.9\%) & Vehicle data links with range, motion, pose,
temperature, or channel sensing & Changing geometry, acquisition, tracking, and
service continuity & \cite{OISAC_SCR00123,OISAC_SCR00273} \\
\addlinespace[1.25pt]
\textbf{Optical access} 32 (15.5\%) & Access and transport links with channel,
event, or fiber state sensing & Reach, synchronization, and traffic continuity &
\cite{OISAC_SCR00007,OISAC_SCR01006,OISAC_SCR01049} \\
\addlinespace[1.25pt]
\textbf{Indoor positioning} 13 (6.3\%) & Indoor optical links with user or robot
position and orientation sensing & Geometry, calibration, coverage, and update
continuity &
\cite{OISAC_SCR00910,OISAC_SCR01024} \\
\addlinespace[1.25pt]
\textbf{Data center} 4 (1.9\%) & Short reach data links with channel estimation
or link state sensing & Traffic continuity and calibrated link references &
\cite{OISAC_SCR00007} \\
\bottomrule
\end{tabularx}

\vspace{3pt}
\begin{tabularx}{\textwidth}{@{}
>{\raggedright\arraybackslash}p{0.16\textwidth}
>{\raggedright\arraybackslash}p{0.39\textwidth}
>{\raggedright\arraybackslash}X
>{\centering\arraybackslash}p{0.095\textwidth}@{}}
\toprule
\multicolumn{4}{@{}l}{\emph{Infrastructure}} \\
\addlinespace[1pt]
\textbf{Domain, n (\%)} & \textbf{Common communication and sensing role} &
\textbf{Key operating condition} & \textbf{Studies} \\
\midrule
\textbf{Environmental monitoring} 55 (26.7\%) & Distributed links with vibration,
temperature, or environmental event sensing & Persistent observation,
localization, and traceable reference data &
\cite{OISAC_SCR00957,OISAC_SCR01006,OISAC_SCR01095} \\
\addlinespace[1.25pt]
\textbf{Smart infrastructure} 41 (19.9\%) & Infrastructure links with structural,
motion, or event monitoring & Continuous observation, calibration, and
resilient alert delivery & \cite{OISAC_SCR00277,OISAC_SCR00763,OISAC_SCR01006} \\
\addlinespace[1.25pt]
\textbf{Industrial} 34 (16.5\%) & Industrial data or control with process,
strain, temperature, motion, or position sensing & Delay, continuity,
calibration, and safe operation & \cite{OISAC_SCR00277,OISAC_SCR00376} \\
\addlinespace[1.25pt]
\textbf{Security and surveillance} 26 (12.6\%) & Monitoring links with anomaly,
event, or location sensing & Coverage, event discrimination, and alert latency &
\cite{OISAC_SCR01006} \\
\bottomrule
\end{tabularx}

\vspace{3pt}
\begin{tabularx}{\textwidth}{@{}
>{\raggedright\arraybackslash}p{0.16\textwidth}
>{\raggedright\arraybackslash}p{0.39\textwidth}
>{\raggedright\arraybackslash}X
>{\centering\arraybackslash}p{0.095\textwidth}@{}}
\toprule
\multicolumn{4}{@{}l}{\emph{Human centered, extreme settings, and residual}} \\
\addlinespace[1pt]
\textbf{Domain, n (\%)} & \textbf{Common communication and sensing role} &
\textbf{Key operating condition} & \textbf{Studies} \\
\midrule
\textbf{Healthcare} 6 (2.9\%) & Wearable or ambient links with vital sign,
activity, or gesture sensing & Motion robustness, physiological reference,
illumination, and user safety & \cite{OISAC_SCR00001,OISAC_SCR00376} \\
\addlinespace[1.25pt]
\textbf{Aerospace} 20 (9.7\%) & Satellite links with atmospheric or wavefront
sensing & Pointing, turbulence, delay, and link adaptation &
\cite{OISAC_SCR00456} \\
\addlinespace[1.25pt]
\textbf{Underwater} 2 (1.0\%) & Underwater link concepts and subsea data links
with temperature or event sensing & Water attenuation or cable calibration,
packaging, and access & \cite{OISAC_SCR00763,OISAC_SCR01095} \\
\addlinespace[1.25pt]
\textbf{Other} 17 (8.3\%) & Emerging links with range, imaging, object, or
channel sensing & Task specific metrics and reporting & \cite{OISAC_SCR00953} \\
\bottomrule
\end{tabularx}
\vspace{1pt}
\parbox{\textwidth}{\raggedright\footnotesize\emph{Note.} Application domains
overlap, so a study may contribute to more than one row. Counts and percentages
describe coverage within the 206 included studies. For each domain, the table
includes a selected study set and lists the corresponding references in the
final column. The 6G and access row identifies the application category, and
the final subsection provides the 6G relevance classification.\par}
\end{table*}

Across the map, the dominant requirement moves with the application. Network
and mobility settings emphasize acquisition, timing, and continuity. Deployed
infrastructure emphasizes persistent observation alongside carried traffic.
Human centered and extreme settings add physiological references, safety,
privacy, pointing, atmosphere, packaging, and access to the test environment.
These requirement bundles show which operating conditions must accompany a
reported result before it can support comparison or reuse.

\subsection{From 6G Relevance to Network Evidence}

The exclusive 6G crosswalk places 138 studies (67.0\%) in the direct class, 64
(31.1\%) in the inferential class, 1 (0.5\%) in the weak class, and 3 (1.5\%)
in the not applicable class. These four labels are review codes for positioning
the corpus. They follow the frozen Phase E crosswalk, which groups source level
codes through explicit or direct and inferred or enabling normalization rules.
Weak and not applicable assignments remain separate. Direct assignments include
\cite{OISAC_SCR00362,OISAC_SCR01049}, and inferential assignments include
\cite{OISAC_SCR00007,OISAC_SCR00647}. The weak and not applicable classes are
documented in
\cite{OISAC_SCR00959,OISAC_SCR00294}.

The distribution connects the reviewed corpus with future access, transport,
localization, automation, environmental awareness, and high frequency carrier
generation. A relevance code therefore locates a study within the 6G
discussion. Standard conformance, interoperability, and deployment readiness
require additional evidence tied to the relevant network framework
\cite{ITUR_M2160_2023}.

Network evidence begins when the optical or photonic subsystem occupies an
explicit end to end role. Access, fronthaul, transport, edge processing, and
service control bring timing, mobility, orchestration, energy, security, and
recovery requirements. A convincing experiment evaluates useful communication
and sensing under the conditions of that role and reports the resources needed
to sustain both. This creates a traceable path from the physical observable and
hardware to the service objective.

Within the scope of this survey, the strongest 6G interpretation follows that
path through both functional outcomes and their operating conditions. It
describes the evidence reported by the 206 included studies. Claims about
industry adoption, economic feasibility, and standard compliance require
evidence designed for those questions. Taken together, this section connects
the physical observable to the application requirement and the network role.
Section~\ref{sec:discussion_roadmap} uses the remaining gaps in that chain to
develop the research and reporting roadmap.
\FloatBarrier
